\documentclass[10pt,a4paper]{article}

\usepackage[margin=1.55cm,top=1.35cm,bottom=1.4cm]{geometry}
\usepackage[T1]{fontenc}
\usepackage[utf8]{inputenc}
\usepackage{lmodern}
\usepackage{microtype}
\usepackage{xcolor}
\usepackage{enumitem}
\usepackage{titlesec}
\usepackage[hidelinks]{hyperref}
\usepackage{bookmark}

\input{glyphtounicode}
\pdfgentounicode=1
\pagestyle{empty}
\setlength{\parindent}{0pt}
\setlength{\parskip}{0pt}
\renewcommand{\familydefault}{\sfdefault}

\definecolor{navy}{HTML}{102F55}
\definecolor{blue}{HTML}{1E5A96}
\definecolor{text}{HTML}{1D2733}
\definecolor{muted}{HTML}{526174}
\definecolor{rule}{HTML}{C9D5E2}

\color{text}
\urlstyle{same}
\hypersetup{
  colorlinks=true,
  urlcolor=blue,
  linkcolor=blue,
  pdftitle={Anwar Ibrahim -- Curriculum Vitae},
  pdfauthor={Anwar Ibrahim},
  pdfsubject={Machine Learning Research and Engineering},
  pdfkeywords={machine learning, reinforcement learning, accelerator control, generative models, agentic AI}
}

\titleformat{\section}
  {\large\bfseries\color{navy}}
  {}{0pt}{}
  [\vspace{2pt}\color{rule}\titlerule]
\titlespacing*{\section}{0pt}{10pt}{6pt}

\setlist[itemize]{
  leftmargin=1.35em,
  itemsep=1.6pt,
  topsep=3pt,
  parsep=0pt,
  partopsep=0pt,
  label=\textcolor{blue}{\small\textbullet}
}
\setlist[enumerate]{
  leftmargin=1.65em,
  itemsep=4pt,
  topsep=3pt,
  parsep=0pt,
  partopsep=0pt
}

\newcommand{\contactsep}{\hspace{0.45em}\textcolor{rule}{\textbullet}\hspace{0.45em}}
\newcommand{\entry}[4]{%
  \textbf{#1}\par
  {\textit{\color{muted}#2\if\relax\detokenize{#3}\relax\else, #3\fi} \textcolor{rule}{\textbar} \textbf{#4}}\par\vspace{-2pt}
}
\newcommand{\project}[3]{%
  \textbf{#1} --- \href{#2}{\textcolor{blue}{GitHub repository}}\\[-1pt]
  {\color{muted}#3}\par
}
\newcommand{\publication}[3]{%
  #1 #2\\[-1pt]
  \textit{#3}
}
\newcommand{\skill}[2]{%
  \textbf{\color{navy}#1:} #2\par\vspace{2pt}
}

\begin{document}

\begin{center}
  {\fontsize{24}{28}\selectfont\bfseries\color{navy} Anwar Ibrahim}\\[3pt]
  {\large\color{blue} Machine Learning Researcher $\mid$ Machine Learning Engineer}\\[6pt]
  {\small
    Moscow, Russia \contactsep \href{mailto:anwar96ibrahim@gmail.com}{anwar96ibrahim@gmail.com}\\[2pt]
    \href{https://anwar9ibrahim.github.io/}{anwar9ibrahim.github.io}
    \contactsep \href{https://github.com/Anwar9Ibrahim}{github.com/Anwar9Ibrahim}
    \contactsep \href{https://www.linkedin.com/in/anwar-ibrahim-7116b4205/}{linkedin.com/in/anwar-ibrahim-7116b4205}
  }\\[2pt]
  {\small\href{https://www.hse.ru/en/org/persons/494515691/}{HSE University profile: hse.ru/en/org/persons/494515691}}
\end{center}

\vspace{-3pt}
\section{Professional Summary}
Machine Learning Researcher and Machine Learning Engineer with experience developing reinforcement learning (RL), research software, production-oriented machine learning, and agentic large language model (LLM) systems. Expertise includes Python, PyTorch, continuous-control methods, simulation environments, benchmarking against classical optimization, model evaluation, deployment, and monitoring. First author of three peer-reviewed journal articles and lecturer in generative models and practical machine learning.

\section{Work Experience}

\entry{Junior Research Fellow}{HSE University --- LAMBDA}{Moscow, Russia}{Oct 2023--Present}
\begin{itemize}
  \item Design and implement reinforcement-learning methods and Python research software for accelerator beamline control.
  \item Build Gymnasium-compatible simulation environments that integrate learning agents with Elegant and Self Describing Data Sets (SDDS) accelerator tools.
  \item Train and benchmark Deep Deterministic Policy Gradient (DDPG), Soft Actor-Critic (SAC), and Twin Delayed DDPG (TD3) agents against differential-evolution and Bayesian-optimization baselines.
  \item Authored three peer-reviewed journal articles as first author on simulation-based accelerator control and present results at international conferences and workshops.
\end{itemize}

\entry{Lecturer --- Generative Models in Machine Learning}{HSE University}{Moscow, Russia}{Jan 2025--Present}
\begin{itemize}
  \item Teach \textit{Generative Models in Machine Learning} in English, covering variational autoencoders, generative adversarial networks (GANs), autoregressive models, and normalizing flows.
  \item Deliver practical machine-learning instruction and contribute to a reinforcement-learning competition at the All-Russian High-Energy Physics School.
\end{itemize}

\entry{AI Engineer --- Agentic Systems}{Eolos AI}{}{Dec 2025--Oct 2026}
\begin{itemize}
  \item Developed components for private agentic artificial intelligence systems built around large language models and multi-step LangGraph workflows.
  \item Designed and implemented human-in-the-loop review and feedback paths, with application data managed in Convex.
  \item Used LangSmith tracing and evaluation to support the observability and reliability of LLM workflows.
\end{itemize}

\entry{Machine Learning Engineer}{EXPentech}{Damascus, Syria}{Jun 2020--Sep 2023}
\begin{itemize}
  \item Developed and evaluated machine-learning models and prepared datasets through exploratory analysis.
  \item Supported model deployment and monitoring in production-oriented machine-learning workflows.
\end{itemize}

\section{Education}

\entry{PhD Candidate, Artificial Intelligence and Machine Learning}{HSE University}{Moscow, Russia}{Nov 2023--Present}
\begin{itemize}
  \item Doctoral research: reinforcement-learning-based control of accelerator beamlines.
\end{itemize}

\entry{Master's in Data Science}{HSE University}{Moscow, Russia}{2021--Jul 2023}
\begin{itemize}
  \item Degree programme in Applied Mathematics and Information Science.
\end{itemize}

\entry{Bachelor's in Computer Software Engineering}{Higher Institute for Applied Sciences and Technology}{Damascus, Syria}{2015--2020}

\newpage
\section{Peer-Reviewed Publications}
\begin{enumerate}
  \item \publication{\textbf{Anwar Ibrahim}, Fedor Ratnikov, Maxim Kaledin, Alexey Petrenko, Denis Derkach.}{\href{https://doi.org/10.1016/j.cpc.2026.110407}{\textit{RL-ABC: Reinforcement Learning for Accelerator Beamline Control}.}}{Computer Physics Communications, 329, 110407 (2026).}

  \item \publication{\textbf{Anwar Ibrahim}, Alexey Petrenko, Maxim Kaledin, Ehab Suleiman, Fedor Ratnikov, Denis Derkach.}{\href{https://doi.org/10.1142/S0218301326410272}{\textit{Reinforcement Learning for Accelerator Beamline Control: A Simulation-Based Approach}.}}{International Journal of Modern Physics E, 35(07) (2026).}

  \item \publication{\textbf{Anwar Ibrahim}, Denis Derkach, Alexey Petrenko, Fedor Ratnikov, Maxim Kaledin.}{\href{https://doi.org/10.1134/S1063779625700716}{\textit{Optimization of the Accelerator Control by Reinforcement Learning: A Simulation-Based Approach}.}}{Physics of Particles and Nuclei, 56(6), 1476--1481 (2025).}
\end{enumerate}

\section{Selected Conference Contributions}
\begin{itemize}
  \item \textbf{CHEP 2026, Bangkok --- Talk:} \href{https://indico.cern.ch/event/1471803/contributions/6966806/}{RL-ABC: Reinforcement Learning for Accelerator Beamline Control}.
  \item \textbf{ACAT 2025 --- Oral contribution:} \href{https://indico.cern.ch/event/1488410/contributions/6562827/}{Optimisation of the Accelerator Control by Reinforcement Learning: A Simulation-Based Approach}.
  \item \textbf{NUCLEUS 2025 --- Talk:} \href{https://indico.spbu.ru/event/1/contributions/171/}{Reinforcement Learning for Accelerator Beamline Control}.
  \item \textbf{CERN Workshop 2025:} \href{https://indico.cern.ch/event/1502120/contributions/6505973/}{Reinforcement Learning for Accelerator Beamline Control}.
  \item \textbf{MMCP 2024:} \href{https://indico.jinr.ru/event/4467/contributions/28919/}{Optimization of Accelerator Control by Reinforcement Learning}.
\end{itemize}

\section{Selected Projects}
\project{RL-ABC}{https://github.com/hse-cs/RL-ABC}{Research framework that converts Elegant accelerator lattice configurations into reinforcement-learning environments, with continuous-control agents and classical optimization baselines.}
\vspace{4pt}
\project{Implicit Q-Learning on Hopper}{https://github.com/Anwar9Ibrahim/offlineRL_hopper}{PyTorch implementation of Implicit Q-Learning on the D4RL Hopper expert dataset, including actor, value, twin-critic, and target networks.}
\vspace{4pt}
\project{Interpretable Reinforcement Learning for Pong}{https://github.com/Anwar9Ibrahim/Interpretable-RL}{Analysis toolkit using saliency maps, Q-value inspection, paddle-position analysis, and an interactive dashboard.}
\vspace{4pt}
\project{Multilingual Contexto}{https://github.com/Anwar9Ibrahim/Contexto-game}{Semantic word-guessing game built with BERT-family models and Streamlit, supporting Arabic, English, and Russian.}

\section{Technical Skills}
\skill{Machine Learning}{Reinforcement learning, generative models, diffusion models, natural language processing (NLP), computer vision, human-in-the-loop systems}
\skill{Programming}{Python, C\#, SQL}
\skill{Frameworks}{PyTorch, TensorFlow, Gymnasium, Stable-Baselines3, Streamlit, OpenCV}
\skill{Agentic AI}{LangChain, LangGraph, LangSmith, Convex, large language model (LLM) workflow orchestration}
\skill{Scientific Tools}{Elegant, SDDS, simulation-based optimization, experiment tracking}
\skill{Engineering}{Git, model evaluation, deployment, monitoring}
\skill{Languages}{Arabic, English, Russian}

\end{document}
